\usepackage{times}
\usepackage{etoolbox}
%\usepackage[left=1in,right=1in,top=1in,bottom=1in,letterpaper]{geometry}
\usepackage[left=1in,right=1in,top=1in,bottom=1.45in,letterpaper]{geometry}
\usepackage{alltt}
\usepackage{amssymb}
\usepackage{amsmath}
\usepackage{amsfonts}
\usepackage{amsthm}
\usepackage{array}
\usepackage{commonunicode}
\usepackage{mathtools}
\usepackage{listings}
% \usepackage[section]{minted}
% \usemintedstyle[julia]{bw} % stata-light
%\setminted{fontsize=\normalsize}
% \ifieeetrans\else
%     \usepackage{appendix}
% \fi
\usepackage{threeparttable}
\usepackage{booktabs}
\usepackage{tabularray} % table without lines
\usepackage{enumitem}
\usepackage{multicol}
\usepackage{ragged2e}

% \ifieeetrans
%     % \usepackage{fontspec}  % requires XeLaTeX or LuaLaTeX
%     % \newfontfamily\headingfont{Georgia}
% \fi
\usepackage{xspace}
\usepackage{fvextra}
\usepackage{graphicx}
\usepackage{soul}
\usepackage{microtype}
\usepackage{fancyhdr}
\usepackage{caption}
\DeclareCaptionLabelSeparator{emdash}{\textemdash}
\usepackage{etoolbox}
\usepackage{parskip}
\usepackage{float}
\usepackage{comment}
\usepackage{colortbl}
\usepackage{csquotes}
\usepackage[svgnames]{xcolor}
% https://c7.alamy.com/comp/3A53X8G/color-palette-26-plum-and-deep-violet-shades-3d-light-soft-to-deep-dark-purple-tones-with-hex-codes-and-names-colour-swatches-scheme-3A53X8G.jpg
\definecolor{DarkWisteria}{RGB}{37,13,49}
\definecolor{BlackViolet}{RGB}{22,15,26}
%\usepackage[hidelinks]{hyperref}
\usepackage{hyperref}
\hypersetup{
    colorlinks=true, % Replaces default red boxes with colored text
    linkcolor=BlackViolet,  % color for internal links (includes Operator names)
    urlcolor=DarkGreen,     % color for external web links
    citecolor=DarkBlue,     % color for bibliography citations
    anchorcolor=Olive
}
\usepackage{pdflscape}
\usepackage{afterpage}
\usepackage{adjustbox}
\usepackage{forloop}
\usepackage{pgffor}
\usepackage{datetime2}
\usepackage{listings}
% \ifieeetrans\else
%     \usepackage{titlesec}
%     \titleformat{\section}
%     {\normalfont\sffamily\Large\bfseries\color{CadetBlue}}
%     {\thesection}{1em}{}
%     \titleformat{\subsection}
%     {\normalfont\sffamily\large\bfseries\color{black}}
%     {\thesubsection}{1em}{}
%     \titleformat{\subsubsection}
%     {\normalfont\sffamily\large\color{black}}
%     {\thesubsubsection}{1em}{}
% \fi

\setcounter{secnumdepth}{8}

\setlength{\parskip}{8pt} % Space between paragraphs
% Be more consistent with use of \val, includign for zero
\usepackage{bbm}

\usepackage{titling}
\pretitle{\begin{flushleft}\Large\bfseries}
        \posttitle{\par\end{flushleft}}
\preauthor{\begin{flushleft}}
        \postauthor{\end{flushleft}}
\predate{\begin{flushleft}}
        \postdate{\end{flushleft}}

%%%%%%%%%%%%%%%%%%%%%%%%%%%%%%%%%%
%% Document-specific macros

% for \shortminus def
\DeclareMathSymbol{\shortminus}{\mathbin}{AMSa}{"39}
\def\sminus{\ensuremath{\mathord{\shortminus}}} % looks better in tables

\def\parameter#1{\ensuremath{\mathsf{#1}}\xspace}
\def\val#1{$\mathsf{#1}$}
\def\code#1{$\mathtt{#1}$}

% \def\Aug#1{#1{\bar{\mathbb R}}}
\def\Aug{\omega}
\def\Cer{\Aug}
\def\zero{0}
\def\one{1}

\def\arctanh{\mathrm{arctanh}}
\def\arcsinh{\mathrm{arcsinh}}
\def\arccosh{\mathrm{arccosh}}

\def\negskip{\vspace*{-\baselineskip}}

\def\fmt{\parameter{fmt}}
\def\emax{\parameter{emax}}
\def\emin{\parameter{emin}}
\def\bias{\parameter{bias}}
\def\SE{\parameter{SE}}
\def\x{$x$\xspace}
\def\s{\sigma} % should be "s", temporarily "S" during typographical edits
\def\p{\parameter{P}} % should be "p", temporarily "P" during typographical edits
\def\w{\parameter{W}} % should be "w", temporarily "W" during typographical edits
\def\t{\parameter{T}} % should be "t", temporarily "T" during typographical edits
\def\k{\parameter{K}} % should be "k", temporarily "K" during typographical edits
\def\b{\parameter{B}} % should be "b", temporarily "B" during typographical edits
\def\minusZero{$\sminus{0}$}

\makeatletter
\def\BinaryK#1#2#3#4{\ensuremath{\mathsf{Binary}\{#1,#2,#3,#4\}}}
\def\binarystring{Binary}
\def\binarypptailp p{\ensuremath{\mathsf{pP}}}
\def\binarypptail#1#2#3{\ensuremath{\mathsf{p#1#2#3}}}
\def\binaryptail p{\@ifnextchar p\binarypptailp\binarypptail}
\def\binarytail{\xspace}
\def\binary#1{\ensuremath{\mathsf{\binarystring#1}}\@ifnextchar p\binaryptail\binarytail}
\def\binaryplain{\ensuremath{\mathsf{\binarystring}}}
\makeatother

\def\IEEEBinaryLong#1#2{\ieeestd\ensuremath{\mathsf{binary#1#2}}}
\def\IEEEBinary#1#2{\ensuremath{\mathsf{binary#1#2}}}
\def\IEEEBinaryK{\ensuremath{\mathsf{binary}\{K\}}}
\def\BFloat16{\ensuremath{\mathsf{BFloat16}}}

\makeatletter
\def\renderop#1{\mathsf{#1}}
\def\runopnormal{}
\def\runopangle<#1>{\textcolor{darkgray}{_{#1}}}
\def\runopnormalQ{}
\def\runopangleQ<#1>{}
% \def\runopnormalA{}
\def\runopangleA<#1>{\langle#1\rangle}
\def\runopangleA<#1>{{\mathtt{<}}#1{\mathtt{>}}}
\def\runoptail{\@ifnextchar<\runopangle\runopnormal}
\def\runoptailQ*{\@ifnextchar<\runopangleQ\runopnormalQ}
\def\runoptailA+{\@ifnextchar<\runopangleA\runopnormalA}
\def\runoptailstar{%
    \@ifnextchar*\runoptailQ%
    {\@ifnextchar+\runoptailA\runoptail}}
\def\runop#1{\hyperlink{sym:#1}{\renderop{#1}}\runoptailstar}
\def\runaugop#1{\hyperlink{sym:#1}{\Aug\renderop{#1}}\runoptailstar}
\makeatother

\def\anchorop#1{\hypertarget{sym:#1}{}}
\newcommand{\anchorops}[1]{\forcsvlist{\anchorop}{#1}}

\newcommand{\declop}[1]{%
    \expandafter\def\csname #1\endcsname{\runop{#1}}}
\newcommand{\declops}[1]{\forcsvlist{\declop}{#1}}

% \newcommand{\newop}[1]{%
%     \anchorop{#1}%
%     \declop{#1}%
% }
% \newcommand{\newops}[1]{\forcsvlist{\newop}{#1}}

% Generic macro to build operation lists
% Creates four macros for managing lists of operations:
%   \<listname> - the raw list with leading comma
%   \declop<singular> - adds a single operation to the list
%   \declop<listname> - adds multiple operations (CSV) to the list
%   \use<listname> - outputs the list without the leading comma
\makeatletter
\newcommand{\defoplist}[2]{% #1 = list name, #2 = singular name
    % Initialize empty list
    \expandafter\def\csname #1\endcsname{}%
    % Define \declop<singular> to add one operation
    \expandafter\newcommand\csname declop#2\endcsname[1]{%
        \declop{##1}%
        % Append to list: preserve existing content, add comma and new operation in math mode
        \expandafter\edef\csname #1\endcsname{%
            \unexpanded\expandafter\expandafter\expandafter{\csname #1\endcsname},
            \noexpand$\expandafter\noexpand\csname ##1\endcsname\noexpand$}%
    }%
    % Define \declop<listname> to add multiple operations from CSV
    \expandafter\newcommand\csname declop#1\endcsname[1]{%
        \forcsvlist{\csname declop#2\endcsname}{##1}}%
    % Define \use<listname> to output list without leading comma
    \expandafter\newcommand\csname use#1\endcsname{%
        \expandafter\expandafter\expandafter\@gobble\csname #1\endcsname}%
}
\makeatother

% Define lists for different kinds of operations
\defoplist{blkunaries}{blkunary}
\defoplist{blkbinaries}{blkbinary}


% start of new opearation declarations

\declop{RoundAway}
\declop{CodeIsOdd}
\declop{FMA}
\declop{FAA}
\declop{Clamp}
\declop{Class}
\declop{reduce}
\declop{ConvertFromBlock}
\declop{ConvertToBlock}
\declop{ConvertToBlockMaxAbsFinite}
\declop{logsumexp}
%
\declops{Extended,Finite,Signed,Unsigned}
%
\declops{True,False}
\declops{IsOdd,IsEven}
\declops{NearestTiesToEven,NearestTiesToAway,TowardPositive,TowardNegative,TowardZero,ToOdd,StochasticA,StochasticB,StochasticC}
\declops{SatFinite,SatPropagate,SatNone}
\declops{PrecisionOf,bitwidthOf,SignednessOf,DomainOf,MaxFiniteOf,MinFiniteOf,MinPositiveOf,MinNormalOf,ExponentBiasOf,ExponentBitsOf,TrailingBitsOf,RoundOf,SatOf}
\declops{CompareLess,CompareLessEqual,CompareEqual,CompareGreaterEqual,CompareGreater}
\declops{IsZero,IsOne,IsInfinite,IsNaN,IsSignMinus,IsNormal,IsSubnormal,IsFinite}
\declops{NextGreaterThan,NextLessThan}
\declops{NextGreaterThanAux,NextLessThanAux}
\declops{BlockReduceAdd,BlockReduceMultiply}
\declops{BlockDotProduct}
%
\declopblkunaries{Convert,Abs,Negate}
\declopblkunaries{Sqrt, RSqrt, Recip}
\declopblkunaries{Exp,Log,ExpMinusOne,LogOnePlus,Exp2,Log2}
\declopblkunaries{Sin,Cos,Tan,ArcSin,ArcCos,ArcTan}
\declopblkunaries{Sinh,Cosh,Tanh,ArcSinh,ArcCosh,ArcTanh}
\declopblkunaries{SinPi,CosPi,TanPi,ArcSinPi,ArcCosPi,ArcTanPi}
\declopblkunaries{Softplus}
%
\declopblkbinary{CopySign}
\declopblkbinaries{Add,Subtract,Multiply,Divide}
\declopblkbinaries{Hypot}
\declopblkbinaries{ArcTan2,ArcTan2Pi}
\declopblkbinaries{Maximum,Minimum,MaximumNumber,MinimumNumber}
\declopblkbinaries{MaximumMagnitude,MinimumMagnitude,MaximumMagnitudeNumber,MinimumMagnitudeNumber}
\declopblkbinaries{MinimumFinite,MaximumFinite}

% end of new operation declarations


\def\ssec#1{{\bf #1}\\[2pt]}
\def\Opnd#1#2{\parbox{12mm}{\hfill {#1} :}~#2}
\def\OpndWide#1#2{{{#1} :}~#2}
%\def\opnd#1:#2;{\Opnd{#1}{#2}}
\def\Case#1{${}\qquad #1$}
%\def\case #1;{\Case{#1}}
\def\commenttext#1{\textcolor{blue}{#1}}
\def\comment#1{\hfill\mbox{\commenttext{\quad ---#1}}}
\def\comment#1{\hfill\mbox{\quad (#1)}}

\def\domain{\Delta}
\def\signedness{\Sigma}

\def\dedent{\hspace*{-1em}}


\def\any{\ast}
\def\sign{\operatorname{sign}}
\def\OneFromBool#1{\mathbbm 1\!\left[#1\right]}
\def\gives{\rightarrow\,\xspace}
\def\bool{\operatorname{Boolean}}

% https://tex.stackexchange.com/questions/83509/hfill-in-math-mode
\makeatletter
\newcommand{\pushright}[1]{\ifmeasuring@#1\else\omit\hfill$\displaystyle#1$\fi\ignorespaces}
\newcommand{\pushleft}[1]{\ifmeasuring@#1\else\omit$\displaystyle#1$\hfill\fi\ignorespaces}
\makeatother

\def\floor#1{\lfloor#1\rfloor}
\def\ceil#1{\lceil#1\rceil}

\def\sumop{Σ}

\def\NInf{\sminus\parameter{Inf}}
\def\Inf{\parameter{Inf}}
\def\NaN{\parameter{NaN}}
\def\NaNs{\parameter{NaNs}}

\makeatletter
\def\NaNinStar*#1{\NaN_{#1}}
\def\NaNinNormal#1{\NaN}
\def\NaNin{\@ifnextchar*\NaNinStar\NaNinNormal}
\makeatother

\newcommand{\newkeyword}[1]{\expandafter\def\csname #1\endcsname{\operatorname{\mathbf{\lowercase{#1}}}}}

\forcsvlist{\newkeyword}{If,Then,Else,End,And,Or,Not,Otherwise}
% \def\Or{\operatorname{\mathbf{or}}}
% \def\And{\operatorname{\mathbf{and}}}
% \def\Not{\operatorname{\mathbf{not}}}

\makeatletter
\def\WhereStar*{\mathbf{where}\quad}
\def\WhereNormal{\quad\WhereStar*}
\def\Where{\@ifnextchar*\WhereStar\WhereNormal}
\makeatother

\newcommand\ieeestdtm{IEEE Std 754\texttrademark\xspace}
\newcommand\ieeestd{IEEE Std 754\xspace}
% \newcommand\ieeenew{IEEE Std 754-2019\xspace}
\newcommand\ieeenew{IEEE Std 754\texttrademark\xspace}

\def\classnormal#1{#1}
\def\classsubnormal#1{\textcolor{RoyalBlue}{#1}}\def\classsubnormalname{blue}
\def\classinfnan#1{\textcolor{Sienna}{#1}}\def\classinfnanname{brown}
\def\classzero#1{\classinfnan{#1}}


% Blue-underlined url in heading font
\def\myurl#1{\href{#1}{\setulcolor{blue}\ul{#1}}}

% Print a float as M \times 2^{E}, with E width fixed to that of "-88"
\def\binaryfloat#1#2{\ensuremath{#1 \times 2^{\mathrlap{#2}\phantom{-88}}}}

\def\binaryfloatinline#1#2{\ensuremath{#1 \times 2^{\mathrlap{#2}}}}


\newcommand{\spec}[1]{{\bf \boldmath  #1}}

% Environment to scope the local \defs
% E.g.
% \begin{scope}
%  \def\stuff{...}
%  use \stuff here
% \end{scope}
\newenvironment{scope}{}{}

% % enable line numbers, if desired
% \usepackage{lineno}
% \linenumbers
% \pagewiselinenumbers

\newcommand{\cmw}[1]{\marginpar{\color{red}#1}}


% IEEE SA draft page style and caption formatting
\setlength{\headheight}{28pt}
\pagestyle{fancy}
\fancyhf{}
\fancyhead[C]{\footnotesize IEEE P3109/D1. March 2026\\Draft Standard for Arithmetic Formats for Machine Learning}
\fancyfoot[C]{\footnotesize \thepage\\Copyright © 2026 IEEE. All rights reserved.\\This is an unapproved IEEE Standards Draft, subject to change.}
\renewcommand{\headrulewidth}{0pt}
\renewcommand{\footrulewidth}{0pt}

\captionsetup[table]{position=above,justification=centering,singlelinecheck=false,labelsep=emdash,textfont=normalfont}
\captionsetup[figure]{position=bottom,justification=centering,singlelinecheck=false,labelsep=emdash,textfont=normalfont}

% Notes formatting
\newcommand{\notetext}[1]{\par\noindent\textbf{NOTE—} #1\par}
\newcommand{\notenumbered}[2]{\par\noindent\textbf{NOTE #1—} #2\par}
\newcommand{\notesfootnote}{\footnote{Notes to text, tables, and figures are for information only and do not contain requirements needed to implement the standard.}}

% Cross-reference helpers
\newcommand{\clauseref}[1]{Clause~\ref{#1}}
\newcommand{\annexref}[1]{Annex~\ref{#1}}

\def\projspec{\rho}

\def\inputaug#1{\input{ieee-p3109/iml2latex/w#1}\negskip}

\def\Project{\runaugop{Project}}
\def\RoundToPrecision{\runaugop{RoundToPrecision}}
\def\Saturate{\runop{Saturate}}
\def\wSaturate{\runaugop{Saturate}}

\def\ConvertWithin{\runop{Convert}}
\def\ConvertIn{\runop{ConvertFromIEEE754}}
\def\ConvertOut{\runop{ConvertToIEEE754}}

%
% IEEE Typography conformance

%   fonts
\usepackage[T1]{fontenc}
\usepackage{tgtermes}
\usepackage{tgheros}

%   customizing
\usepackage{titlesec}
\usepackage{caption}
\usepackage{tocloft}

%   headings
\titleformat*{\section}{\sffamily\fontsize{12pt}{14pt}\selectfont\bfseries}
\titleformat*{\subsection}{\sffamily\fontsize{11pt}{13pt}\selectfont\bfseries}
%% !!FIXME!! (Andrew have you done some redefining of \subsubsection?)
%\titleformat*{\subsubsection}{\sffamily\fontsize{10pt}{12pt}\selectfont\bfseries}

% Caption setup
\captionsetup{font=normalsize} % 10pt in a 10pt document
